\begin{frame}{Opbrækning af vildt}
    \begin{itemize}
        \item \textbf{Formål}: Fjerne indvolde hurtigt for at undgå gæring, forurening og kødsurhed
        \item Gør det \textbf{straks} efter nedlæggelse (helst inden 30–60 min.)
        \item Grundtrin (klovbærende vildt):
          \begin{itemize}
            \item Læg dyret på ryggen (ben opad) – gerne på rent underlag
            \item Skær fra brystben til hale (åbn bughulen)
            \item Undgå at skære i tarmene/vommen – skær forsigtigt rundt om endetarm
            \item Udtag indvolde samlet (hjerte, lunger, lever, mave, tarme)
            \item Fjern spiserør og luftrør op mod halsen
            \item Tør bughulen ren med klud/papir (ingen vand!)
          \end{itemize}
        \item \textbf{Hygiejne}: Rene hænder/kniv, undgå jord/græs-kontakt
        \item Efter opbrækning: Afkøl hurtigt – hæng op i skygge/køligt sted
        \item https://www.facebook.com/reel/889956353611580
    \end{itemize}
\end{frame}



